% FILE: 00_project_identity.tex
% Project: lifecycle
% Role: process-lifecycle-state-definitions-and-transitions

\section{Project Identity}
\label{sec:lifecycle-identity}

\subsection{Purpose}
\label{sec:lifecycle-identity-purpose}

The \texttt{lifecycle} project defines the complete, formally-grounded process intelligence
lifecycle as a directed state machine over typed stage tokens. Its purpose is to specify
every admissible state a process model may occupy---from initial M\&A acquisition ingestion
through cryptographic final archival---and to govern every transition between those states
using six quality gates derived from the Blue River Dam doctrine. The framework is the
authority layer that determines when a process model may advance, must halt, or must be
refused, and it establishes the mathematical criteria that distinguish a board-defensible
claim from an unverifiable narrative.

The framework is not aspirational architecture. Every state, transition rule, gate criterion,
and taxonomy documented in the \texttt{lifecycle/} directory was authored between 31 May and
1 June 2026, with 43 Markdown source files covering every stage and supporting taxonomy. The
purpose of this explicit scope is to make the lifecycle surface auditable: an independent
reviewer can enumerate the files and verify that no stage, no transition, and no taxonomy
is absent.

\subsection{Repository Surface}
\label{sec:lifecycle-identity-surface}

The lifecycle directory resides at \texttt{\textasciitilde/process-intelligence/lifecycle/}
and contains 43 files, all in Markdown format. The surface is organized into three layers.
The first layer covers the six core operational stages: Design, Simulation, Monitoring,
Repair, Optimization, and Decommissioning. Each stage has a canonical definition file
(e.g., \texttt{define\_design-state\_process\_intelligence.md}) and a companion narrative
file (e.g., \texttt{design.md}). The second layer covers the seven supporting states:
Acquisition, Construction, Activation, Integration, Operation, Archive, and Board-Projection,
each defined in its own file. The third layer covers the five taxonomies---Process Asset,
Process Debt, Process Residual, Process Risk, and False-Claim---plus the autonomic control
documents: \texttt{MAPE\_K\_MAP.md}, \texttt{define\_blue\_river\_dam\_lifecycle\_gate\_map.md},
\texttt{define\_autonomic\_knowledge\_actuation\_map.md}, and the maturity model
\texttt{MATURITY\_MODEL.md}.

No TTL/RDF ontology files, no executable test files, and no cryptographic receipt JSON
files exist within the \texttt{lifecycle/} subdirectory itself. Receipt JSON artifacts
reside in the parent \texttt{receipts/} directory, and the parent-level
\texttt{checkpoints/PROCESS\_INTELLIGENCE\_ALIVE\_001.md} issues the program-level verdict.

\subsection{Primary Research Role}
\label{sec:lifecycle-identity-role}

The lifecycle project occupies the role of \emph{process-lifecycle-state-definitions-and-transitions}
within the dissertation. It answers the question: what are the lawful states of a process
model under process intelligence governance, and under what formal criteria may a model
transition from one state to the next? This question is not decorative. Without a formally
specified lifecycle, the wasm4pm execution engine has no authority surface against which
to enforce state transitions, and the MAPE-K autonomic loop has no admissible stage
boundaries to respect. The lifecycle specification is, in this sense, the type system of
the entire process intelligence architecture.

The lifecycle project is also the primary surface on which the dissertation's M\&A due
diligence argument is grounded. The ACQUISITION.md document defines the diligence gap
(the inability of most sellers to replay their own process claims), and the lifecycle
framework provides the staged artifact requirements---model evidence, log evidence,
conformance evidence, repair evidence, and replay verification---that close that gap.

\subsection{Key Artifacts}
\label{sec:lifecycle-identity-artifacts}

The key artifacts produced by the lifecycle project are as follows. First, the evidence
lifecycle typed state machine (v30.1.1 Spec), expressed as a directed graph over the set
$\mathcal{S} = \{\texttt{Raw}, \texttt{Parsed}, \texttt{Admitted}, \texttt{Refused},
\texttt{Projected}, \texttt{Exportable}, \texttt{Receipted}\}$ with lawful transitions
enforced at the type level. Second, the six Blue River Dam quality gates, which govern
stage transitions from Design through Decommissioning. Third, the MAPE-K mapping matrix,
which assigns each of the six operational stages to its corresponding autonomic loop
component. Fourth, the five-level Process Intelligence Maturity Model (L1 Ad-hoc through
L5 Autonomic) and the six-level Process Readiness Levels (PRL 0--5). Fifth, the Process
Debt quantification formula $D_p = \sum_{c \in L}(C_{\mathrm{deviate}}(c) +
C_{\mathrm{delay}}(c) + C_{\mathrm{overhead}}(c))$ and its M\&A valuation implications.
Sixth, the Cryptographic Decommissioning Receipt specification $R_d =
\mathrm{Ed25519}_{K_{\mathrm{priv}}}(\mathrm{BLAKE3}(N) \parallel \mathrm{BLAKE3}(L_{\mathrm{final}})
\parallel C_{\mathrm{total}} \parallel F_{\mathrm{final}} \parallel T_{\mathrm{retire}})$.
Seventh, the LIFECYCLE\_RECEIPT in the parent \texttt{receipts/RECEIPT\_REGISTRY.md},
certifying that all 41 or more lifecycle phase files were produced with correct input/output
type coverage.
